\documentclass[UTF8,11pt]{ctexart}
\usepackage[a4paper,margin=24mm]{geometry}
\usepackage{amsmath,amssymb,amsthm,mathtools,mathrsfs}
\usepackage[all]{xy}
\usepackage{xurl,hyperref,enumitem}
\usepackage{tikz}
\usetikzlibrary{arrows.meta,cd,decorations.markings,calc,patterns}
\hypersetup{hidelinks,breaklinks=true}
\setlength{\emergencystretch}{3em}
\allowdisplaybreaks
\newtheorem*{theorem}{Theorem}
\newtheorem*{thm}{Theorem}
\newtheorem*{lemma}{Lemma}
\newtheorem*{proposition}{Proposition}
\newtheorem*{prop}{Proposition}
\newtheorem*{corollary}{Corollary}
\newtheorem*{cor}{Corollary}
\newtheorem*{claim}{Claim}
\theoremstyle{definition}
\newtheorem*{definition}{Definition}
\newtheorem*{defn}{Definition}
\newtheorem*{remark}{Remark}
\newtheorem*{example}{Example}
\newcommand{\R}{\mathbb R}
\newcommand{\C}{\mathbb C}
\newcommand{\Z}{\mathbb Z}
\newcommand{\Q}{\mathbb Q}
\newcommand{\N}{\mathbb N}
\newcommand{\F}{\mathbb F}
\newcommand{\D}{\mathbb D}
\newcommand{\bB}{\mathbb B}
\newcommand{\bC}{\mathbb C}
\newcommand{\bR}{\mathbb R}
\newcommand{\bZ}{\mathbb Z}
\newcommand{\bN}{\mathbb N}
\newcommand{\bQ}{\mathbb Q}
\newcommand{\bD}{\mathbb D}
\newcommand{\bF}{\mathbb F}
\newcommand{\CP}{\mathbb{CP}}
\newcommand{\RP}{\mathbb{RP}}
\newcommand{\sO}{\mathscr O}
\newcommand{\sI}{\mathscr I}
\newcommand{\sF}{\mathscr F}
\newcommand{\sG}{\mathscr G}
\newcommand{\sH}{\mathscr H}
\newcommand{\sR}{\mathscr R}
\newcommand{\sM}{\mathscr M}
\newcommand{\sV}{\mathscr V}
\newcommand{\sU}{\mathscr U}
\DeclareMathOperator{\Hom}{Hom}
\DeclareMathOperator{\Ext}{Ext}
\DeclareMathOperator{\Tor}{Tor}
\DeclareMathOperator{\Spec}{Spec}
\DeclareMathOperator{\Proj}{Proj}
\DeclareMathOperator{\supp}{supp}
\DeclareMathOperator{\im}{im}
\DeclareMathOperator{\id}{id}
\DeclareMathOperator{\Tr}{Tr}
\DeclareMathOperator{\tr}{tr}
\DeclareMathOperator{\rank}{rank}
\DeclareMathOperator{\ord}{ord}
\DeclareMathOperator{\dist}{dist}
\DeclareMathOperator{\End}{End}
\DeclareMathOperator{\Aut}{Aut}
\DeclareMathOperator{\Gal}{Gal}
\DeclareMathOperator{\vol}{vol}
\DeclareMathOperator{\Cl}{Cl}
\DeclareMathOperator{\coker}{coker}
\DeclareMathOperator{\codim}{codim}
\DeclareMathOperator{\divergence}{div}
\DeclareMathOperator{\Ric}{Ric}
\DeclareMathOperator{\grad}{grad}
\DeclareMathOperator{\Hess}{Hess}
\newcommand{\abs}[1]{\left\lvert#1\right\rvert}
\newcommand{\sabs}[1]{\lvert#1\rvert}
\newcommand{\babs}[1]{\bigl\lvert#1\bigr\rvert}
\newcommand{\Babs}[1]{\Bigl\lvert#1\Bigr\rvert}
\newcommand{\bbabs}[1]{\biggl\lvert#1\biggr\rvert}
\newcommand{\BBabs}[1]{\Biggl\lvert#1\Biggr\rvert}
\newcommand{\norm}[1]{\left\lVert#1\right\rVert}
\newcommand{\snorm}[1]{\lVert#1\rVert}
\newcommand{\bnorm}[1]{\bigl\lVert#1\bigr\rVert}
\newcommand{\Bnorm}[1]{\Bigl\lVert#1\Bigr\rVert}
\newcommand{\bbnorm}[1]{\biggl\lVert#1\biggr\rVert}
\newcommand{\BBnorm}[1]{\Biggl\lVert#1\Biggr\rVert}
\newcommand{\widebar}[1]{\overline{#1}}
\newcommand{\nicefrac}[2]{\frac{#1}{#2}}
\newcommand{\linnprod}[2]{\langle#1,#2\rangle}
\newcommand{\myindex}[1]{#1}
\newcommand{\glsadd}[1]{}
\newcommand{\avoidbreak}{}
\newcommand{\sourceref}[1]{\textnormal{[原文：\nolinkurl{#1}]}}
\newcommand{\thmref}[1]{\sourceref{#1}}
\newcommand{\lemmaref}[1]{\sourceref{#1}}
\newcommand{\propref}[1]{\sourceref{#1}}
\newcommand{\corref}[1]{\sourceref{#1}}
\newcommand{\defnref}[1]{\sourceref{#1}}
\newcommand{\exampleref}[1]{\sourceref{#1}}
\newcommand{\exerciseref}[1]{\sourceref{#1}}
\newcommand{\figureref}[1]{\sourceref{#1}}
\newcommand{\sectionref}[1]{\sourceref{#1}}
\newcommand{\appendixref}[1]{\sourceref{#1}}
\newcommand{\stacksref}[2]{\href{https://stacks.math.columbia.edu/tag/#1}{#2 [#1]}}


\newcommand{\E}{\mathbb E}
\newcommand{\Prob}{\mathbb P}
\newcommand{\Reff}{\mathcal R_{\mathrm{eff}}}
\newcommand{\eps}{\varepsilon}
\DeclareMathOperator{\diam}{diam}
\DeclareMathOperator{\Var}{Var}
\DeclareMathOperator{\Crit}{Crit}
\newcommand{\dd}{\,\mathrm d}

\begin{document}
\section*{062\quad 中心单代数嵌入的Skolem--Noether共轭定理}
\subsection*{来源与计数目标}
The Stacks Project Authors, \emph{The Stacks Project}, Chapter 11, Sections 11.4 and 11.6. 主命题 Theorem 11.6.1 (074Q)，附 Lemmas 11.4.6 (074E)、11.4.7 (074F) 的全文与证明。实际访问在线版本：2026-09-28。\par\url{https://stacks.math.columbia.edu/tag/074E}\par\url{https://stacks.math.columbia.edu/tag/074F}
\par\url{https://stacks.math.columbia.edu/tag/074Q}
\par\textbf{本文件只计一个问题：}两个单代数到同一有限维中心单代数的 $k$-代数嵌入相差内共轭；不是只证明矩阵代数自同构的特殊情形。
\subsection*{作者命题与人类证明}
\begin{lemma}[11.4.6, Tag 074E]
Let $A$ be a finite simple $k$-algebra.
\begin{enumerate}
\item There exists exactly one simple $A$-module $M$ up to isomorphism.
\item Any finite $A$-module is a direct sum of copies of a simple module.
\item Two finite $A$-modules are isomorphic if and only if they have the same dimension over $k$.
\item If $A=\operatorname{Mat}(n\times n,K)$ with $K$ a finite skew field extension of $k$, then $M=K^{\oplus n}$ is a simple $A$-module and $\operatorname{End}_A(M)=K^{op}$.
\item If $M$ is a simple $A$-module, then $L=\operatorname{End}_A(M)$ is a skew field finite over $k$ acting on the left on $M$, we have $A=\operatorname{End}_L(M)$, and the centers of $A$ and $L$ agree. Also $[A:k][L:k]=\dim_k(M)^2$.
\item For a finite $A$-module $N$ the algebra $B=\operatorname{End}_A(N)$ is a matrix algebra over the skew field $L$ of (5). Moreover $\operatorname{End}_B(N)=A$.
\end{enumerate}
\end{lemma}
\begin{proof}
By Theorem 11.3.3 we can write $A=\operatorname{Mat}(n\times n,K)$ for some finite skew field extension $K$ of $k$. By Lemma 11.4.5 the category of modules over $A$ is equivalent to the category of modules over $K$. Thus (1), (2), and (3) hold because every module over $K$ is free. Part (4) holds because the equivalence transforms the $K$-module $K$ to $M=K^{\oplus n}$. Using $M=K^{\oplus n}$ in (5) we see that $L=K^{op}$.

The statement about the center of $L=K^{op}$ follows from Lemma 11.4.5. The statement about $\operatorname{End}_L(M)$ follows from the explicit form of $M$. The formula of dimensions is clear. Part (6) follows as $N$ is isomorphic to a direct sum of copies of a simple module.
\end{proof}

\begin{lemma}[11.4.7, Tag 074F]
Let $A$, $A'$ be two simple $k$-algebras one of which is finite and central over $k$. Then $A\otimes_k A'$ is simple.
\end{lemma}
\begin{proof}
Suppose that $A'$ is finite and central over $k$. Write $A'=\operatorname{Mat}(n\times n,K')$, see Theorem 11.3.3. Then the center of $K'$ is $k$ and we conclude that $A\otimes_k K'$ is simple by Lemma 11.4.4. Hence $A\otimes_k A'=\operatorname{Mat}(n\times n,A\otimes_k K')$ is simple by Lemma 11.4.5.
\end{proof}

\begin{theorem}[11.6.1, Skolem--Noether, Tag 074Q]
Let $A$ be a finite central simple $k$-algebra. Let $B$ be a simple $k$-algebra. Let $f,g:B\to A$ be two $k$-algebra homomorphisms. Then there exists an invertible element $x\in A$ such that $f(b)=xg(b)x^{-1}$ for all $b\in B$.
\end{theorem}
\begin{proof}
Choose a simple $A$-module $M$. Set $L=\operatorname{End}_A(M)$. Then $L$ is a skew field with center $k$ which acts on the left on $M$, see Lemmas 11.3.2 and 11.4.6. Then $M$ has two $B\otimes_k L^{op}$-module structures defined by $m\cdot_1(b\otimes l)=lmf(b)$ and $m\cdot_2(b\otimes l)=lmg(b)$. The $k$-algebra $B\otimes_k L^{op}$ is simple by Lemma 11.4.7. Since $B$ is simple, the existence of a $k$-algebra homomorphism $B\to A$ implies that $B$ is finite.

Thus $B\otimes_k L^{op}$ is finite simple and we conclude the two $B\otimes_k L^{op}$-module structures on $M$ are isomorphic by Lemma 11.4.6. Hence we find $\varphi:M\to M$ intertwining these operations. In particular $\varphi$ is in the commutant of $L$ which implies that $\varphi$ is multiplication by some $x\in A$, see Lemma 11.4.6. Working out the definitions we see that $x$ is a solution to our problem.
\end{proof}
\subsection*{整理说明（非作者正文）}
保留作者先修支撑结果与主证明，不把两个支撑引理另计题数。允许的标准先修为 Wedderburn 单代数结构定理（11.3.3）、矩阵代数的 Morita 等价（11.4.5）、Schur 引理（11.3.2）以及中心斜域张量积的双边理想性质（11.4.4）。11.4.4 的具体内容是：若 $K$ 是中心为 $k$ 的斜域，则 $A\otimes_kK$ 的双边理想均为 $J\otimes_kK$，其中 $J$ 是 $A$ 的双边理想。定位：\url{https://stacks.math.columbia.edu/tag/074C}。

难度依据：在上述结构理论之外，本题仍需把两个嵌入编码为同一向量空间上的两种张量积模结构，以等维模的同构构造交织算子，再通过双中心化子把它还原成原代数中的可逆共轭元。不是直接引用 Skolem--Noether 或其等价结论。

主证明的左右作用书写、$L^{op}$ 与 $lmf(b)$ 均按作者正文保留；网页评论不混入正文。外部编号保留并给出可定位 Tag；这里只取得网页数学正文，保存的副本是整理转录，不冒称整章原生文件。
\subsection*{可修改的简短标注（非作者正文）}
$c$：单代数嵌入的共轭分类。

$a$：张量积模；对立代数；简单模分类；Morita 等价；交织算子；双中心化子；内共轭。

\texttt{cross\_theory\_status = yes}。将嵌入的比较转成张量积代数上的模同构，再把模交织算子识别为原中心单代数中的可逆元。位置：Theorem 11.6.1 两段证明及 11.4.6(3)(5)。
标签为非穷尽、可修改初稿。
\subsection*{许可与修改记录}
原作者及作品见来源栏；来源采用 GNU Free Documentation License 1.2 或后续版本。本题证摘录和附加整理沿用该许可。2026-09-28：增加题号、来源定位、独立排版、整理说明和初稿标注；数学内容的取材范围及排版变化见整理说明。
\end{document}
